%! Modeling with Exponential Functions: Interest, Half-Life, and e — Unit 7, Lesson 7.4 — Exit Ticket
\documentclass[10pt]{article}
\usepackage{algebra2-article}
\usepackage{algebra2-boxes}
\newcommand{\TallMath}[1]{$\displaystyle #1\rule[-1.4em]{0pt}{3.2em}$}

\begin{document}

\pageheader{Unit 7, Lesson 7.4}{Exit Ticket}
\namedateperiod

\vspace{0.08in}

No notes. On items 1 and 2, write the substitution before you write the answer --- the model \emph{is}
part of the answer. Round money to the nearest cent.

\begin{enumerate}[leftmargin=1.8em, itemsep=11pt]
  \item \textbf{Compound interest.} \$4000 is invested at $3\%$ compounded \textbf{monthly} for
        $5$ years.

        \vspace{4pt}
        Base $1+\frac{r}{n}=$ \blank{2.6cm} \qquad Exponent $nt=$ \blank{1.8cm}

        \vspace{5pt}
        $A=$ \blank{5.0cm} $=$ \blank{2.4cm}

  \item \textbf{Half-life.} A $50$ mg dose of a drug has a half-life of $4$ days.

        \vspace{4pt}
        $A(t)=$ \blank{4.2cm} \hfill Amount left after $12$ days: \blank{2.2cm}

        \vspace{5pt}
        Will the amount in the patient's body ever reach exactly $0$ mg according to this model?
        Circle one --- \textbf{yes} \; / \; \textbf{no} --- and name the feature of the graph that
        settles it. \writeline

  \item \textbf{Multiple choice.} An investment of \$1500 earns $4.8\%$ annual interest compounded
        \textbf{quarterly}. Which expression gives its value after $7$ years?
        \begin{enumerate}[label=\Alph*., leftmargin=1.9em, itemsep=1pt]
          \item $1500(1.048)^{7}$
          \item $1500(1.012)^{7}$
          \item $1500(1.048)^{28}$
          \item $1500(1.012)^{28}$
        \end{enumerate}

  \item \textbf{The wall.} To find when that \$1500 doubles, you would solve
        $1500(1.012)^{4t}=3000$, which becomes $(1.012)^{4t}=2$. Yesterday's common-base method cannot
        finish this. Does the equation still \emph{have} an answer? Circle one --- \textbf{yes} \; /
        \; \textbf{no} --- and say how you know. \writeline
\end{enumerate}

\end{document}
